\documentclass[12pt]{article}
\usepackage[utf8]{inputenc}
\usepackage[T1]{fontenc}
\usepackage[margin=1in]{geometry}
\usepackage{setspace}
\usepackage{xcolor}
\usepackage[colorlinks=true, urlcolor=blue!60!black, linkcolor=blue!60!black]{hyperref}
\setlength{\parindent}{2em}
\setlength{\parskip}{1.25em}
\renewcommand{\baselinestretch}{1.0}
\begin{document}
\begin{center}
{\LARGE \textit{United Public Workers v. Mitchell}, 330 U.S. 75 (1947)\par}
\vspace{0.8em}
{\large Case Brief --- Four-Page Edition\par}
\end{center}
\vspace{1em}

\section{Memory Jogger}

A 4--3 Court dismissed eleven federal employees' pre-enforcement challenges as nonjusticiable and upheld the Hatch Act's ban on partisan political management as applied to a U.S. Mint worker, declaring that Ninth and Tenth Amendment objections fail whenever Congress exercises a granted power (\href{https://supreme.justia.com/cases/federal/us/330/75/}{Justia}).

\section{Detailed Case Facts}

The Hatch Political Activity Act of 1939, as amended in 1940, barred most executive-branch employees from taking ``any active part in political management or in political campaigns,'' while preserving their right to vote and to express opinions on political subjects. Violation required removal. Twelve federal employees and the United Public Workers of America sued the Civil Service Commission, claiming violations of the First Amendment, of political rights reserved by the Ninth and Tenth Amendments, and of Fifth Amendment due process.

Eleven plaintiffs, ranging from statisticians to a lens grinder and a welder, had \emph{not} violated the Act; they wished, off duty, to write political articles, ring doorbells, and serve as poll watchers, but refrained for fear of dismissal. Only George P. Poole, a roller at the U.S. Mint in Philadelphia, had violated it: he served as a Democratic ward executive committeeman, worked the polls on election day in 1940, and paid other party workers. In 1944 the Commission charged him and proposed his removal; Poole admitted the charges.

\section{Procedural History}

A three-judge district court in the District of Columbia found all claims justiciable but upheld \S~9(a) and granted summary judgment for the Commission. The appeal was docketed late, and the Government challenged the Court's jurisdiction. After argument and reargument, the Supreme Court affirmed on February 10, 1947.

\section{Judicial Votes}

Affirmed, \textbf{4--3} on the merits; \textbf{5--2} against hearing the eleven pre-enforcement claims.
\begin{itemize}
\item \textbf{Majority:} Justice Reed, joined by Chief Justice Vinson and Justices Burton and Frankfurter.
\item \textbf{Concurrence:} Justice Frankfurter, on jurisdiction.
\item \textbf{Dissents:} Justice Black; Justice Douglas (in part); Justice Rutledge (as to Poole).
\item \textbf{Not participating:} Justices Murphy and Jackson.
\end{itemize}

\section{Holding}

(1) The 60-day docketing provision is not jurisdictional. (2) Employees who merely desire to engage in prohibited activity, without violating the Act or facing a concrete threat, present no Article III case or controversy. (3) Poole, who admitted violations and faced a proposed removal, presented a justiciable controversy. (4) Congress may, within reasonable limits, bar executive employees, including off-duty industrial workers, from active partisan management to promote an efficient, disciplined civil service, and ``[i]f granted power is found,'' Ninth and Tenth Amendment objections ``must fail.'' Reed's concession that Congress could not bar a ``Republican, Jew or Negro'' from federal office is dicta.

\section{Analysis of Opinions}

\textbf{Majority (Reed).} Courts may not issue advisory opinions on an ``abstract'' attack; they require ``definite rights'' and ``definite prejudicial interferences.'' Poole's case was concrete because removal was mandatory on his admitted facts. On the merits, the Ninth and Tenth Amendments mirror enumerated powers: federal powers are ``subtracted from the totality of sovereignty originally in the states and the people.'' Building on \textit{Ex parte Curtis} and \textit{United States v. Wurzbach}, the Act extended bans on ``political contributions of money'' to ``political contributions of energy.'' Because employees could still vote and speak on public affairs, Congress could rationally fear that even an industrial worker might be swept into a political machine. Treating industrial workers differently was a ``matter of detail for Congress.''

\textbf{Frankfurter (concurring).} The docketing deadline was a congressional limit on appellate jurisdiction the Court could not enlarge, but ``under compulsion of the Court's assumption of jurisdiction'' he joined on the merits.

\textbf{Black (dissenting).} The Act muzzled about three million federal employees, and even their spouses, making them ``mere spectators.'' Restrictions on speech must be narrowly drawn to a ``grave and imminent danger''; Congress could punish coercion directly.

\textbf{Douglas (dissenting in part).} The eleven faced an immediate threat of discharge. A Mint roller is as remote from policymaking ``as a charwoman,'' and federal employees ``are not second-class citizens.''

\textbf{Rutledge (dissenting as to Poole)} agreed with Black but would not reach the other eleven claims.

\section{Examples: Future Applications}

\begin{enumerate}
\item \textbf{Same-Side (Upheld): Partisan Candidacy.} A federal customs inspector is barred from running as a party nominee for the state legislature. Under \textit{Mitchell} and \textit{Letter Carriers} (1973), the ban is valid.
\item \textbf{Same-Side (Nonjusticiable): Vague Desire.} A postal clerk who wishes to join ``unspecified future political campaigns,'' with no planned conduct or warning, presents no case or controversy.
\item \textbf{Opposite-Side (Invalid): Religion or Party Bans.} A rule barring employees from attending Mass or belonging to a mainstream party falls within Reed's own concession.
\item \textbf{Opposite-Side (Invalid): Honoraria Ban.} A ban on pay for off-duty, nonpartisan speeches and articles fails under \textit{United States v. NTEU} (1995).
\item \textbf{Fence-Sitter: Issue-Advocacy Blog.} A park mechanic's blog urges votes against incumbents without endorsing a party. Reed protected expression ``not an objective of party action,'' but electoral advocacy may look like energy directed toward party success.
\end{enumerate}

\section{Critique}

\textbf{Originalist critique.} Randy E. Barnett argues that Reed's ``rights-powers conception'' erases the Ninth Amendment, which protects natural rights that constrain enumerated powers. Kurt T. Lash contends that Reed broke with a tradition, rooted in Madison and Story, treating the Ninth Amendment as a rule of construction limiting federal power. \textit{Griswold v. Connecticut} (1965) later invoked the Ninth Amendment to protect marital privacy.

\textbf{Progressive critique.} \textit{Mitchell} is faulted as a relic of the right-privilege distinction of Holmes's \textit{McAuliffe} epigram, applying near-rational-basis deference to core political speech without narrow tailoring. Later cases dismantled that premise, and Congress relaxed the Act in 1993.

\section{Key Quotations}

\begin{quote}
``The powers granted by the Constitution to the Federal Government are subtracted from the totality of sovereignty originally in the states and the people. Therefore, when objection is made that the exercise of a federal power infringes upon rights reserved by the Ninth and Tenth Amendments, the inquiry must be directed toward the granted power under which the action of the Union was taken. If granted power is found, necessarily the objection of invasion of those rights, reserved by the Ninth and Tenth Amendments, must fail.''
\end{quote}
\noindent --- Justice Reed, for the Court, 330 U.S. at 95--96.

\begin{quote}
``The determination of the extent to which political activities of governmental employees shall be regulated lies primarily with Congress. Courts will interfere only when such regulation passes beyond the generally existing conception of governmental power.''
\end{quote}
\noindent --- Justice Reed, for the Court, 330 U.S. at 102.

\begin{quote}
``The section of the Act here held valid reduces the constitutionally protected liberty of several million citizens to less than a shadow of its substance. It relegates millions of federal, state, and municipal employees to the role of mere spectators of events upon which hinge the safety and welfare of all the people, including public employees.''
\end{quote}
\noindent --- Justice Black, dissenting, 330 U.S. at 115.

\end{document}
